\section{Audited Markov instruments and lawful edits}\label{sec:model}

\subsection{The native transition law}

Fix nonempty finite sets \(M,X,R\) of marks, base states and register cells.
The source state is \((g,u)\in X\times R\). Fix a finite family
\(\chi_i:X\to\{-1,1\}\), indexed by \(I\), with
\[
 \sum_{g\in X}\chi_i(g)=0,\qquad
 \sum_{g\in X}\chi_i(g)\chi_j(g)=|X|\,\mathbf 1_{i=j}.
\]
The nontrivial characters of a finite binary abelian group give examples,
and the proof uses only the displayed identities. Fix a mark reversal
\(m\mapsto\bar m\) of order at most two and a full-support rational
register law \(\nu\). A law consists of rational mark rates \(\theta_m\),
responses \(\alpha_{mi}\), and conditional register tables \(L_m(g,h;u,v)\).
Fix rational \(\delta>0\) and \(0\leq\kappa<1\).
Require
\begin{equation}\label{eq:guards}
 \theta_m\ge\delta,\quad \sum_m\theta_m=1,\qquad
 \sum_i|\alpha_{mi}|\le\kappa,\quad
 \alpha_{\bar m,i}=\alpha_{mi}.
\end{equation}
Write
\begin{equation}\label{eq:native}
 Q_m(g,h)=\frac{1+\sum_i\alpha_{mi}\chi_i(g)\chi_i(h)}{|X|},
 \qquad
 P_m((g,u),(h,v))=\theta_m Q_m(g,h)L_m(g,h;u,v).
\end{equation}
The tables \(L_m(g,h)\) have nonnegative entries, unit row sums and
stationary margin \(\nu\). Their reciprocal pairing is
\[
 \nu(u)L_m(g,h;u,v)=\nu(v)L_{\bar m}(h,g;v,u).
\]
There is one such table for each \(m,g,h\), so the register dynamics may
depend on the base transition.

The positive uniform term in \(Q_m\) has a probabilistic implementation.
Refresh the base uniformly with probability \(1-\sum_i|\alpha_{mi}|\).
With probability \(|\alpha_{mi}|\), choose a uniform endpoint with
\(\chi_i(h)=\sgn(\alpha_{mi})\chi_i(g)\), omitting zero-weight choices.
Balance gives exactly \eqref{eq:native}. In particular
\(Q_m(g,h)\ge(1-\kappa)/|X|>0\).
The unmarked kernel obtained by summing the marked entries is stationary
under \(\mu(g,u)=\nu(u)/|X|\). Indeed the conditional stationary margins
sum out the incoming register, after which the symmetric base kernel
preserves the uniform base law. The reciprocal identity is
\[
 \mu(g,u)P_m((g,u),(h,v))\,\theta_{\bar m}
 =
 \mu(h,v)P_{\bar m}((h,v),(g,u))\,\theta_m .
\]
This fixes the marked reversal relation even when the rates of reversed marks
differ.

\subsection{Held tests and their compiled equations}\label{sec:tests}

A timed word \(w\) is a nonempty finite sequence of marks of length at most
a fixed horizon. Let \(n_m(w)\) be the multiplicity of \(m\) in \(w\).
Its probability from a prepared state is
\(\pi_\theta(w)=\prod_m\theta_m^{n_m(w)}\).
After summation over the final register, its measured probability with
projected endpoint \(h\), from a preparation \((g,u)\), is
\begin{equation}\label{eq:word-return}
 \pi_\theta(w)\,(Q_{w_1}\cdots Q_{w_{|w|}})(g,h).
\end{equation}
This follows by iterative use of the register row sums. The base matrices
share the uniform projection and the orthogonal rank-one projections onto
the probes \(\chi_i\). Multiplying the base matrices multiplies the
eigenvalue on each probe, giving
\[
 \lambda_i(w)=\prod_m\alpha_{mi}^{n_m(w)}.
\]
These coordinates determine the word matrix, since the uniform term is
fixed and the other projections vanish on the orthogonal complement.

A \emph{license} names a comparison of two timed words of equal length.
The held comparison asks for equal word probabilities and equal conditional
projected endpoint laws. By \eqref{eq:word-return} it is also equivalent
to equality of the summed native word probabilities for all preparations
and all projected endpoints. Summing over endpoints recovers each
positive word mass, so the equality can be divided by that mass.
A fixed policy on subsets of the finite license catalogue specifies which
sets may be held.

The distinction between current and continuation tests also arises in
the study of predictive quotients for fixed-support route
transport~\citep{routeholonomy}. Here the licenses require specified
pairs of timed experiments to agree as the transition law changes.

A \emph{publication} is a named query with a stated rational value
that can be checked on the law. The catalogue consists of positive
word-mass ratios for equal-duration pairs and calibrations of the
form \(\alpha_{mi}=c\). For the latter, from any prepared native state
\((g,u)\),
\begin{equation}\label{eq:signal}
 \chi_i(g)\sum_{h,v} P_m((g,u),(h,v))\chi_i(h)=\theta_m\alpha_{mi}.
\end{equation}
This is a signed measurement. Dividing by the observed mark mass verifies the
calibration. Thus the held numerical tests on the rates are positive
homogeneous monomial equations, since the two degrees in a word comparison
equal the two lengths. The spectral tests are signed monomial equations on
\(\alpha\). They include the reversal equalities in \eqref{eq:guards}.
These two systems share the catalogue names and have separate numerical
variables.

\subsection{Capabilities and the edit relation}\label{sec:capabilities}

A profile \(p=(S,F,\Lambda,\mathcal A)\) records the positive native
proposal support \(S\), a schedule \(F=(F_0,\ldots,F_n)\) of partitions
of \(R\), the held licenses \(\Lambda\), and the publications \(\mathcal A\).
A decision at grade \(j\) is a function \(R\to\{0,1\}\) constant on the
classes of \(F_j\). As grades advance, the partitions may only refine.
For each available decision \(d\), require
\begin{equation}\label{eq:decision-mass}
 \sum_{v:\,d(v)=1} L_m(g,h;u,v)=\sum_{v:\,d(v)=1}\nu(v)
 \qquad\text{for all }m,g,h,u .
\end{equation}
These conditional tests are measured from the prepared transition table
because \(\theta_m Q_m(g,h)>0\). Equivalently it suffices to check
the terminal partition. Its block masses add to the masses of every
coarser disclosed block.

A state \(A\) is a profile \(p_A\) with a law that passes its tests,
has exactly its support, and respects the license policy. A profile is
\emph{attained} when some lawful rational state has that profile.
For profiles use the order
\begin{equation}\label{eq:profile-order}
 p\Refines q\quad\Longleftrightarrow\quad
 S_q\subseteq S_p,\quad
 F_j(q)\text{ refines }F_j(p)\text{ for all }j,\quad
 \Lambda_p\subseteq\Lambda_q,\quad \mathcal A_p\subseteq\mathcal A_q .
\end{equation}
The support is observable as the positive entries of \(P\). The
partitions are also observable from the available decisions.
Two cells lie in the same class precisely when every available
decision gives them the same value. Indicators of classes distinguish
different classes.

For a positive rational tolerance \(\varepsilon\), define elementary
edits on lawful states without using role labels. At fixed profile,
a local edit satisfies
\[
 \Norm{P_A-P_B}
 =\max_{m,s,t}|P^A_m(s,t)-P^B_m(s,t)|<\varepsilon .
\]
The other edits follow a cover of the finite order of attained profiles.
A cover is a strict increase with no attained profile strictly between
its endpoints. If the support is unchanged, the prepared transition
table stays fixed. If the support contracts, the projected marked
probabilities \(\theta_m Q_m(g,h)\) stay fixed. The conditional register may
change in that discrete edit. A chain means finite composition of
elementary edits, including the empty chain. We write \(A\leadsto_\varepsilon B\)
and call \(B\) \emph{factorable from \(A\)} if \(A\leadsto_\varepsilon B\)
for every positive rational \(\varepsilon\).

This definition keeps the numerical tolerance on the prepared law,
where it can be tested directly. Every step is directed in the
profile order, so an inherited license or publication remains held
at each later state.

\begin{theorem}[Audited Markov reachability]\label{thm:native}
Fix an instrument as above and let \(A,B\) be lawful rational states. Let
\(\SCore_{p_A}\) denote the signed core of \Cref{def:signed-core} for the
spectral equations compiled from the initial profile \(p_A\), including
the fixed reversal tests. Then
\[
 \bigl(\forall \varepsilon\in\QQ_{>0},\ A\leadsto_\varepsilon B\bigr)
 \quad\Longleftrightarrow\quad
 p_A\Refines p_B\quad\text{and}\quad
 \SCore_{p_A}(\alpha_A)=\SCore_{p_A}(\alpha_B).
\]
The criterion can be recovered from the measured experiments.
\end{theorem}

The next section gives the signed-core construction and proves its
rational completeness. Section~\ref{sec:native} proves the transfer
to native laws and capability changes.
